\section{Extensión estocástica}
En esta sección se analizará el problema desde un punto de vista más amplio, ahora en lugar de tener decisiones y  costes fijos en un escenario determinista, ahora esto variara dependiendo del escenario, ya que en lugar de un solo escenario se tienen 3 en el que se varia la demanda, estando el escenario de demanda normal, el de demanda baja y el de demanda alta.

\subsection{Decisiones de primera etapa}
Las decisiones de primera etapa son aquellas que son independientes del escenario a enfrentar, osea, son decisiones que se toman antes de partir con la distribucion de combustibles, las cuales se formulan de la siguiente forma:

\begin{itemize}
    \item $y_k \in \{0,1\}$: vale 1 si el camión $k \in K$ es utilizado y 0 en caso contrario.

    \item $x_{ijk} \in \{0,1\}$: vale 1 si el camión $k \in K$ viaja directamente desde el nodo $i \in V$ al nodo $j \in V$, y 0 en caso contrario.

    \item $w_{jk} \in \{0,1\}$: vale 1 si la estación $j \in N$ es visitada por el camión $k \in K$, y 0 en caso contrario.

    \item $z_{kcp} \in \{0,1\}$: vale 1 si el compartimento $c \in C$ del camión $k \in K$ es asignado para transportar el producto $p \in P$, y 0 en caso contrario.

    \item $L_{kc0} \geq 0$: cantidad inicial de combustible cargada en el compartimento $c \in C$ del camión $k \in K$ al momento de salir desde el depósito.

    \item $t_{jk} \geq 0$: hora programada de llegada del camión $k \in K$ a la estación $j \in N$.

    \item $u_{jk} \geq 0$: variable auxiliar utilizada para la eliminación de subtours mediante la formulación MTZ.

    \item $s_o, f_o \geq 0$: tiempos de inicio y término de la operación de carga $o \in O$ en el depósito.

    \item $C_{\max} \geq 0$: tiempo de finalización del proceso de carga en el depósito (\textit{makespan}).

    \item $y_{oo'} \in \{0,1\}$: variable de precedencia que indica el orden de utilización de las bahías de carga entre las operaciones $o$ y $o'$.
\end{itemize}

\subsection{Decisiones de segunda etapa}
Las decisiones de segunda etapa son las que se toman y ajustan una vez se conoce el escenario de demanda ${s} \in {S}$, esto permite ajustar la planificación inicial y hace que el modelo pueda reaccionar mediante las variables de recourse, con el fin de absorber la incertidumbre por medio de la penalización de demanda no satisfecha (shortage).

Las variables que dependen del escenario ${s}$ se definen como:

\begin{itemize}
    \item $q_{kcpjs} \geq 0$: cantidad de producto $p \in P$ entregada desde el compartimento $c \in C$ del camión $k \in K$ a la estación $j \in N$ bajo el escenario $s \in {S}$.

    \item $s_{pjs} \geq 0$: demanda no satisfecha (\textit{shortage}) del producto $p \in P$ en la estación $j \in N$ bajo el escenario $s \in {S}$.

    \item $L_{kcjs} \geq 0$: carga remanente en el compartimento $c \in C$ del camión $k \in K$ inmediatamente después de visitar la estación $j \in N$ bajo el escenario $s \in {S}$.
\end{itemize}

\subsection{Función objetivo en extensión estocastica}

La función objetivo de la extensión estocástica a diferencia de la determinista, radica en que, en el modelo determinista, la función objetivo esta dada por costes fijos a consecuencia del unico escenario en el que esta planteada y en el modelo estocástico, la función objetivo esta dividida por dos etapas. La primera de ellas viene dada por los costes deterministas(coste fijo por camiones y distancia recorrida) y la segunda etapa se presenta como el valor esperado de las penalizaciones de segunda etapa (coste por demanda insatisfecha) a lo largo de cada escenario ponderado por sus respectivas probabilidades.

La función objetivo estocástica se presenta de la siguiente manera:

$$\min Z = \sum_{k \in K} F_k y_k + c_d \sum_{k \in K} \sum_{i \in V} \sum_{j \in V, j \neq i} d_{ij} x_{ijk} + \mathbb{E}_s \left[ Q(x, y, z, L_{\cdot\cdot 0}, s) \right]$$

Al expandir el término de valor esperado usando las probabilidades de los escenarios $p_s$, la función objetivo se formula como:

$$\min Z = \sum_{k \in K} F_k y_k + c_d \sum_{k \in K} \sum_{i \in V} \sum_{j \in V, j \neq i} d_{ij} x_{ijk} + \sum_{s \in S} p_s \left( c_s \sum_{j \in N} \sum_{p \in P} s_{jps} \right)$$

Donde:
\begin{itemize}
    \item $\sum_{k \in K} F_k y_k$: Costo fijo total de los camiones utilizados (Primera Etapa).
    
    \item $c_d \sum_{k \in K} \sum_{i \in V} \sum_{j \in V, j \neq i} d_{ij} x_{ijk}$: Costo total por distancia recorrida por la flota (Primera Etapa).
    
    \item $\sum_{s \in S} p_s \left( c_s \sum_{j \in N} \sum_{p \in P} s_{jps} \right)$: Costo esperado de penalización por shortage de combustible no entregado (Segunda Etapa).
\end{itemize}

\subsection{Clasificación de restricciones por etapa}

\subsubsection{Restricciones de primera etapa}
Las restricciones de primera etapa son aquellas que deben cumplirse de manera estricta en la planificacion y no contienen variables que dependen de los escenarios.

\begin{itemize}
    \item \textbf{Conservación de flujo en ruteo:}
    \[
    \sum_{j \in N} x_{0jk} = y_k,
    \qquad \forall k \in K
    \]
    \[
    \sum_{i \in N} x_{i0k} = y_k,
    \qquad \forall k \in K
    \]
    \[
    \sum_{i \in V,\, i \neq j} x_{ijk} = w_{jk},
    \qquad \forall j \in N,\ \forall k \in K
    \]
    \[
    \sum_{h \in V,\, h \neq j} x_{jhk} = w_{jk},
    \qquad \forall j \in N,\ \forall k \in K
    \]

    \item \textbf{Unicidad de visita a las estaciones:}
    \[
    \sum_{k \in K} w_{jk} \leq 1,
    \qquad \forall j \in N
    \]

    \item \textbf{Eliminación de subtours (MTZ):}
    \[
    u_{ik} - u_{jk} + |N|x_{ijk} \leq |N| - 1,
    \qquad \forall k \in K,\ \forall i,j \in N,\ i \neq j
    \]
    \[
    1 \leq u_{jk} \leq |N|w_{jk},
    \qquad \forall j \in N,\ \forall k \in K
    \]

    \item \textbf{Ventanas de tiempo y arribo:}
    \[
    t_jk \geq a_j - M(1-w_{jk}),
    \qquad \forall j \in N,\ \forall k \in K
    \]
    \[
    t_jk \leq b_j + M(1-w_{jk}),
    \qquad \forall j \in N,\ \forall k \in K
    \]
    \[
    t_jk \geq t_ik + T_i + \frac{d_{ij}}{v_k} - M(1-x_{ijk}),
    \qquad \forall k \in K,\ \forall i,j \in N,\ i \neq j
    \]
    \[
    t_jk \geq h_k + \frac{d_{0j}}{v_k} - M(1-x_{0jk}),
    \qquad \forall j \in N,\ \forall k \in K
    \]

    \item \textbf{Compatibilidad compartimento-producto en carga:}
    \[
    \sum_{p \in P} z_{kcp} \leq 1,
    \qquad \forall k \in K,\ \forall c \in C
    \]

    \item \textbf{Capacidad de compartimentos para carga inicial:}
    \[
    L_{kc0} \leq Q_{kc}\sum_{p \in P} z_{kcp},
    \qquad \forall k \in K,\ \forall c \in C
    \]

    \item \textbf{Estabilidad de carga inicial en el depósito (antes de salir):}
    \[
    \frac{L_{kc0}}{Q_{kc}} - \frac{L_{kc'0}}{Q_{kc'}} \leq \Delta,
    \qquad \forall k \in K,\ \forall c,c' \in C,\ c \neq c'
    \]
    \[
    \frac{L_{kc'0}}{Q_{kc'}} - \frac{L_{kc0}}{Q_{kc}} \leq \Delta,
    \qquad \forall k \in K,\ \forall c,c' \in C,\ c \neq c'
    \]

    \item \textbf{Scheduling de Carga en el Depósito:} Todas las restricciones asociadas al makespan $(C_{\max})$, tiempos de limpieza, secuencia interna de camiones y no solapamiento en bahías.
\end{itemize}

\subsubsection{Restricciones de segunda etapa}
Las restricciones de segunda etapa deben ser cumplidas para cada uno de los escenarios y vinculan las decisiones de ruteo o carga inicial con la demanda que es realizada en cada uno de los escenarios.

\begin{itemize}
    \item \textbf{Satisfacción de demanda con shortage:}
    \[
    \sum_{k \in K}\sum_{c \in C} q_{kcpsj} + S_{jps} = D_{jps},
    \qquad \forall j \in N,\ \forall p \in P,\ \forall s \in S
    \]

    \item \textbf{Compatibilidad física de entregas:}
    \[
    q_{kcpsj} \leq Mz_{kcp},
    \qquad \forall k \in K,\ \forall c \in C,\ \forall p \in P,\ \forall j \in N,\ \forall s \in S
    \]
    \[
    q_{kcpsj} \leq Mw_{jk},
    \qquad \forall k \in K,\ \forall c \in C,\ \forall p \in P,\ \forall j \in N,\ \forall s \in S
    \]

    \item \textbf{Acoplamiento de carga (el total entregado en ruta no puede superar lo cargado en el depósito):}
    \[
    \sum_{j \in N}\sum_{p \in P} q_{kcpsj} \leq L_{kc0},
    \qquad \forall k \in K,\ \forall c \in C,\ \forall s \in S
    \]

    \item \textbf{Conservación y actualización de carga en ruta por escenario:}
    \[
    L_{kcjs} \leq L_{kcis} - \sum_{p \in P} q_{kcpsj} + M(1-x_{ijk}),
    \qquad \forall k \in K,\ \forall c \in C,\ \forall i \in V,\ \forall j \in N,\ i \neq j,\ \forall s \in S
    \]
    \[
    L_{kcjs} \geq L_{kcis} - \sum_{p \in P} q_{kcpsj} - M(1-x_{ijk}),
    \qquad \forall k \in K,\ \forall c \in C,\ \forall i \in V,\ \forall j \in N,\ i \neq j,\ \forall s \in S
    \]

    \item \textbf{Estabilidad de carga en ruta por escenario:}
    \[
    \frac{L_{kcjs}}{Q_{kc}} - \frac{L_{kc'js}}{Q_{kc'}} 
    \leq \Delta + M(1-w_{jk}),
    \qquad \forall k \in K,\ \forall c,c' \in C,\ c \neq c',\ \forall j \in N,\ \forall s \in S
    \]
    \[
    \frac{L_{kc'js}}{Q_{kc'}} - \frac{L_{kcjs}}{Q_{kc}} 
    \leq \Delta + M(1-w_{jk}),
    \qquad \forall k \in K,\ \forall c,c' \in C,\ c \neq c',\ \forall j \in N,\ \forall s \in S
    \]
\end{itemize}

